\documentclass[11pt,a4paper]{article}

\input{style.tex}

\begin{document}

% --- Title Page ---
\begin{titlepage}
    \centering
    \vspace*{2cm}
    {\huge\bfseries\sffamily OE44185: Numerical Methods for Offshore and Dredging Engineering\par}
    \vspace{0.75cm}
    {\Large\sffamily Project Report Template\par}
    \vspace{1.25cm}
    {\large\itshape Modelling a Floating Offshore Wind Turbine (OC3 Phase-IV Spar-Buoy)\par}
    \vspace{2cm}
    \begin{tabular}{ll}
        \textbf{Project title:} & [Insert project title] \\
        \textbf{Group number:} & [Insert group number] \\
        \textbf{Student names:} & [Insert names and student numbers] \\
        \textbf{Submission date:} & [Insert date] \\
    \end{tabular}
    \vfill
    {\large\sffamily TU Delft -- Offshore Engineering\par}
\end{titlepage}

\pagenumbering{roman}
\setmaxreportpages{15}

\section*{How to Use This Template}
\instructionplaceholder{Template use}{Keep the main report within 15 pages, move long derivations and raw outputs to appendices, and ensure each claim is supported by numerical evidence, verification, or literature comparison.}

\instructionplaceholder{Project scope}{This template is focused on the numerical modelling workflow of OE44185. Your report should cover structural modelling, numerical method derivation, implementation details, and quantitative result analysis.}

\subsection*{Core + Choice Project Architecture}
All groups must complete the \textbf{Core Model} and then select \textbf{at least two Choice Packages}.
\begin{itemize}
    \item Core Model (mandatory): 2D FOWT structural model, baseline wave/wind/current forcing, FEM/FVM formulation, dynamic response assessment and reduced-order/modal response extension.
    \item Choice Packages (select at least two): nonlinear mooring model, irregular-wave loading, time-dependent thrust model, circular-cylinder CFD extension (either for floater or mooring lines), or local mesh refinement strategy.
\end{itemize}

\subsection*{Section Boundary Rules}
\instructionplaceholder{What belongs where}{Keep evidence in the correct section: assumptions and scope in Section 1; discrete numerical formulation in Section 2; implementation and algorithmic choices in Section 3; verification (convergence, stability, modal) before interpreted engineering conclusions in Section 4. Results presented without prior verification are not considered final evidence.}

\subsection*{Coding Ownership and AI-Minimisation}
\instructionplaceholder{Required practice}{Document your own implementation decisions. For every major solver block, state: (i) why this method was selected, (ii) one alternative that was rejected, and (iii) one verification check used to validate behaviour. This is required even when AI tools are used for language support or debugging assistance.}

\subsection*{Written Statement: Restricted Use of Generative AI}
The use of Generative AI tools (for example Copilot, ChatGPT, and Gemini) is restricted in this course report for the following activities:
\begin{itemize}
    \item Generating derivations, final model equations, or numerical conclusions without independent understanding and verification.
    \item Producing full report sections that are meant to evidence your own reasoning and technical judgement.
    \item Producing complete code or solver routines for submission without testing, validating, and understanding the implementation.
\end{itemize}
Students remain fully responsible for all submitted content. If AI is used, complete the in-depth acknowledgement section at the end of this report.

\tableofcontents
\newpage

\pagenumbering{arabic}

\section*{Executive Summary}
\instructionplaceholder{Executive summary evidence}{Write this last. In at most one page, summarise model scope, key numerical choices, validation outcomes, dominant dynamic response conclusions, and main limitations.}

% ============================================================
% Section 1: Model Definition
% Project focus: OC3 phase-IV spar-buoy floating wind turbine
% ============================================================

\section{Model Definition}
\pagecheck{3--4}

\subsection{Assumptions and scope boundaries}
\instructionplaceholder{Required evidence}{List physical assumptions and scope limits before presenting equations or implementation choices. State what is included, what is excluded, expected impact of each simplification, and the operational range for which your model is intended.}

\subsection{Geometry and system definition}
\instructionplaceholder{Required evidence}{Define the 2D model geometry and all major components: RNA point mass, tower model, spar-buoy body, and mooring representation. State dimensions, reference coordinates, and parameter sources used for the OC3/NREL 5MW representation.}

\subsection{Structural model components}
\instructionplaceholder{Pass-level requirements}{Include at minimum: RNA point mass, thin-beam tower, rigid-body spar-buoy, and spring-damper mooring representation calibrated from quasi-static analysis.}

\instructionplaceholder{Choice packages}{Describe optional extensions selected by your team, such as coupled axial-bending tower behaviour, geometrically nonlinear mooring strings, or a flexible-beam floater representation. Justify each extension physically and numerically, and explain why it was prioritised over alternatives.}

\subsection{Environmental loading model}
\instructionplaceholder{Wave, wind, and current forcing}{State the loading models used and where they are applied in the structural model. Include baseline models (Airy waves, Morison loading, equivalent nacelle thrust, mean current force from CFD) and any higher-fidelity options (irregular waves, time-dependent thrust/current loading).}

\subsection{Governing equations and model scope}
\instructionplaceholder{Required evidence}{Present governing equations, state variables, coupling interfaces, and unknowns for the coupled system. Keep this section at the continuous-model level only; do not include discretisation details here.}

% ============================================================
% Section 2: Numerical Discretization
% ============================================================

\section{Numerical Discretization}
\pagecheck{3--4}

\instructionplaceholder{Derivation overview}{Derive the discrete formulation for both subsystems at a level that demonstrates method ownership: Finite Element Method for structural dynamics and Finite Volume Method for current loading. Keep this section focused on mathematical discretization, not software implementation.}

\subsection{Finite Element formulation for the structure}
\instructionplaceholder{Required evidence}{Describe discretisation of structural domains, interpolation functions, weak form, element matrices, boundary conditions, and external force mapping. Justify element types and approximation spaces for the expected dynamic range.}

\subsection{Finite Volume formulation for current loading}
\instructionplaceholder{Required evidence}{Define control-volume discretisation, flux treatment, pressure-velocity coupling, boundary conditions, and force extraction procedure on the cylinder. Include how pressure and viscous contributions are computed.}

\subsection{Coupling strategy and assumptions}
\instructionplaceholder{Required evidence}{Explain how current-force outputs inform structural loading. Clarify one-way or two-way coupling, data transfer, update frequency, and consistency assumptions. Discuss expected implications on accuracy and computational cost.}

\subsection{Discretization justification and alternatives}
\instructionplaceholder{Required evidence}{For each major discretization choice, state why it was selected, one credible alternative, and the expected trade-off in accuracy, robustness, or computational cost.}

% ============================================================
% Section 3: Numerical Schemes
% ============================================================

\section{Numerical Schemes and Implementation}
\pagecheck{2--3}

\instructionplaceholder{Implementation logic}{Document how the discrete equations are implemented in Python or MATLAB. The reader should be able to understand solver workflow, data flow, and time integration choices from this section alone.}

\subsection{Time integration and solver settings}
\instructionplaceholder{Required evidence}{State time-stepping methods, stability constraints, solver tolerances, and convergence criteria. Explain why the selected scheme is appropriate for the coupled structural and loading dynamics.}

\subsection{Algorithm design and code structure}
\instructionplaceholder{Required evidence}{Describe the computational pipeline, key functions/modules, and reproducibility setup. Highlight efficiency choices and any extensions beyond lecture notebooks.}

\subsection{Implementation ownership checklist}
\instructionplaceholder{Required evidence}{For each major solver component, report: selected approach, one alternative considered, and why the final choice was implemented. Keep detailed convergence, stability, and modal verification in Section 4.}

% ============================================================
% Section 4: Numerical Verification and Results Interpretation
% ============================================================

\section{Numerical Verification and Results Interpretation}
\pagecheck{4--5}

\instructionplaceholder{Result strategy}{Present evidence in two stages: (i) numerical verification (convergence, stability, and modal checks), then (ii) interpreted engineering results under relevant operating conditions. Do not present final engineering conclusions before verification evidence is shown.}

\subsection{Numerical verification}
\instructionplaceholder{Required evidence}{Summarise verification outcomes for the implemented model: consistency checks, benchmark comparisons, and error trends that justify the credibility of the final simulations.}

\subsection{Convergence and time-step assessment}
\instructionplaceholder{Required evidence}{Provide mesh and time-step sensitivity analysis. Quantify numerical uncertainty and justify selected discretisation levels for reported final results.}

\subsection{Stability and modal behaviour}
\instructionplaceholder{Required evidence}{Report numerical stability behaviour and natural frequencies/mode shapes of the structural model. Relate dominant modes to expected forced-response characteristics under wave, wind, and current loading.}

\subsection{Interpreted structural response and validation}
\instructionplaceholder{Required evidence}{Compare key structural responses against literature or trusted references. Include displacement/motion measures at relevant locations (for example nacelle and center of gravity) and discuss agreement, deviations, and engineering relevance.}

\subsection{Current-loading analysis}
\instructionplaceholder{Required evidence}{Present drag coefficient trends with inlet velocity, transient shedding behaviour, and Strouhal number analysis where relevant. Compare with literature values and identify likely causes of mismatch.}

\subsection{Limitations and improvements}
\instructionplaceholder{Required evidence}{State model limitations and propose concrete improvements to geometry, coupling fidelity, boundary conditions, or numerical schemes for future work.}


\section*{References}
\instructionplaceholder{Reference practice}{Include only sources that support your model formulation, numerical choices, validation strategy, and interpretation of results. Prioritise course material, standards, and peer-reviewed literature.}

\newpage
\appendix

\section{Appendix A: Additional Derivations and Model Data}
\instructionplaceholder{Appendix use}{Place expanded derivations, parameter tables, and supplementary equations here when they are too long for the main report body.}

\section{Appendix B: Reproducibility Material}
\instructionplaceholder{Appendix use}{Provide enough details to reproduce your results: solver settings, mesh/time-step tables, script structure, and links to executable files or notebooks.}

\section{Appendix C: Decision Log and Verification Checklist}
\instructionplaceholder{Appendix use}{Track key modelling decisions and their verification evidence. Use this log to show technical ownership and reduce black-box coding practices.}

\begin{table}[h!]
    \centering
    \caption{Decision log template.}
    \begin{tabularx}{\textwidth}{>{\raggedright\arraybackslash}p{0.20\textwidth} >{\raggedright\arraybackslash}X >{\raggedright\arraybackslash}p{0.22\textwidth} >{\raggedright\arraybackslash}p{0.22\textwidth}}
        \toprule
        Decision area & Selected approach and rationale & Alternative considered and rejected & Verification evidence (test/plot/table) \\
        \midrule
        Structural model & [Fill in] & [Fill in] & [Fill in] \\
        CFD model & [Fill in] & [Fill in] & [Fill in] \\
        Coupling strategy & [Fill in] & [Fill in] & [Fill in] \\
        Time integration & [Fill in] & [Fill in] & [Fill in] \\
        \bottomrule
    \end{tabularx}
\end{table}

\section*{Acknowledgement of AI Use}
\instructionplaceholder{AI declaration policy}{Use this section whenever AI use is restricted or prohibited for any graded deliverable. Place it at the end of the report (or as an attachment). State clearly which tools were used, for which purposes, and how outputs were checked and adapted. You remain fully responsible for all submitted content.}

\textbf{AI use acknowledgement (mandatory)}
\begin{itemize}
    \item I acknowledge the use of [tool name and link] to support the development of this assessment by: [Select one or more: Explaining concepts, Refining language, Assisting with programming and debugging, Acting as a brainstorming partner, Providing feedback, Other (please specify)].
    \item Specific use: [Describe exactly where AI was used in this report, for example section drafting support, code debugging, equation checks, language editing, or structuring arguments.]
    \item Resulting output: [State the concrete outputs produced with AI support, for example Table X, Figure Y caption revision, code block in Appendix B, or rewritten paragraph in Section Z.]
    \item Verification and adaptation: [Explain what was checked, corrected, or replaced by the authors before submission, including references, calculations, and engineering judgement applied.]
\end{itemize}

\end{document}
